% matvec_tree.tex: how Gila applies G to p. Each axis splits every branch in
% two (bit b = 0: plain, b = 1: half-phase twiddle phi), giving 8 length-N
% branches that are multiplied by their own diagonal symbol and merged back.
% Uses the palette and styles defined in preamble.tex.
\begin{tikzpicture}[font=\small,
    nd/.style={circle, fill=inkcol, inner sep=0pt, minimum size=1.5mm},
    br0/.style={inkcol, line width=0.6pt},
    br1/.style={accent, line width=0.9pt},
    mul/.style={draw=vaccol, fill=vaccol!8, rounded corners=2pt, inner sep=2pt,
      font=\footnotesize, minimum height=4.4mm},
    pathstep/.style={circle, draw=inkcol, fill=white, font=\scriptsize, inner sep=0.6pt,
      minimum size=3.6mm},
  ]
  \def\dy{0.72}
  % column positions
  \def\Xa{0} \def\Xb{1.9} \def\Xc{3.8} \def\Xd{5.7}
  \def\Xe{8.9} \def\Xf{10.8} \def\Xg{12.7} \def\Xh{14.6}
  % node heights: level 1 (b_x), level 2 (b_x b_y), level 3 (b_x b_y b_z)
  \pgfmathsetmacro\top{3.5*\dy}

  % ---- one path through the tree, (b_x b_y b_z) = (011), numbered in the key ----
  \draw[black!12, line width=7pt, line cap=round, line join=round]
    (\Xa,0) -- (\Xb,2*\dy) -- (\Xc,\dy) -- (\Xd,0.5*\dy) -- (\Xe,0.5*\dy)
    -- (\Xf,\dy) -- (\Xg,2*\dy) -- (\Xh,0);

  % ---- forward: split + FFT along x, then y, then z ----
  \foreach \i in {0,1} {
    \pgfmathsetmacro\yi{(2-4*\i)*\dy}
    \draw[br\i] (\Xa,0) -- (\Xb,\yi);
    \draw[br\i] (\Xg,\yi) -- (\Xh,0);
    \foreach \j in {0,1} {
      \pgfmathsetmacro\yj{(3-2*(2*\i+\j))*\dy}
      \draw[br\j] (\Xb,\yi) -- (\Xc,\yj);
      \draw[br\j] (\Xf,\yj) -- (\Xg,\yi);
      \foreach \k in {0,1} {
        \pgfmathtruncatemacro\leaf{4*\i+2*\j+\k}
        \pgfmathsetmacro\yk{(3.5-\leaf)*\dy}
        \draw[br\k] (\Xc,\yj) -- (\Xd,\yk);
        \draw[br\k] (\Xe,\yk) -- (\Xf,\yj);
        \draw[vaccol, line width=0.6pt] (\Xd,\yk) -- (\Xe,\yk);
        \node[mul] at ({(\Xd+\Xe)/2},\yk) {$\tilde{\bm{\mathsf{G}}}_\circ^{(\i\j\k)}$};
        \node[nd] at (\Xd,\yk) {}; \node[nd] at (\Xe,\yk) {};
      }
      \node[nd] at (\Xc,\yj) {}; \node[nd] at (\Xf,\yj) {};
    }
    \node[nd] at (\Xb,\yi) {}; \node[nd] at (\Xg,\yi) {};
  }
  \node[nd] at (\Xa,0) {}; \node[nd] at (\Xh,0) {};
  \node[left, inner sep=3pt] at (\Xa,0) {$\vb*{\mathsf{p}}$};
  \foreach \n/\x/\y in {1/\Xb/2, 2/\Xc/1, 3/\Xd/0.5, 4/\Xe/0.5, 5/\Xf/1, 6/\Xg/2} {
    \node[pathstep] at (\x-0.3,\y*\dy+0.3) {\n};
  }
  \node[right, inner sep=3pt] at (\Xh,0) {$\vb*{\mathsf{f}}$};

  % ---- stage headers: what each edge applies ----
  \pgfmathsetmacro\hy{\top+0.75}
  \foreach \xa/\xb/\op in {
      \Xa/\Xb/{\bm{\mathsf{F}}_x\phi_x^{b_x}},
      \Xb/\Xc/{\bm{\mathsf{F}}_y\phi_y^{b_y}},
      \Xc/\Xd/{\bm{\mathsf{F}}_z\phi_z^{b_z}},
      \Xe/\Xf/{\tfrac12\phi_z^{-b_z}\bm{\mathsf{F}}_z^{-1}},
      \Xf/\Xg/{\tfrac12\phi_y^{-b_y}\bm{\mathsf{F}}_y^{-1}},
      \Xg/\Xh/{\tfrac12\phi_x^{-b_x}\bm{\mathsf{F}}_x^{-1}}} {
    \node[font=\footnotesize] at ({(\xa+\xb)/2},\hy) {$\op$};
  }
  \node[font=\footnotesize] at ({(\Xd+\Xe)/2},\hy) {multiply};

  % ---- Gila functions doing each half ----
  \pgfmathsetmacro\by{-\top-0.45}
  \draw[decorate, decoration={brace, mirror, amplitude=4pt}, inkcol]
    (\Xa,\by) -- (\Xd,\by) node[midway, below=5pt, font=\footnotesize]
    {\texttt{sptBrn!} + FFT};
  \draw[decorate, decoration={brace, mirror, amplitude=4pt}, inkcol]
    (\Xd+0.2,\by) -- (\Xe-0.2,\by) node[midway, below=5pt, font=\footnotesize]
    {\texttt{mulBrn!}};
  \draw[decorate, decoration={brace, mirror, amplitude=4pt}, inkcol]
    (\Xe,\by) -- (\Xh,\by) node[midway, below=5pt, font=\footnotesize]
    {inverse FFT + \texttt{mrgBrn!}};

  % ---- key for the numbered path ----
  \newcommand\stepref[1]{\tikz[baseline=(c.base)] \node[pathstep] (c) {#1};}
  \node[anchor=north west, font=\footnotesize] at (\Xa-0.3,\by-0.85) {%
    \begin{tabular}{@{}c@{\enspace}l@{}}
      \stepref{1} & $\mathsf{p}^{(0)} = \bm{\mathsf{F}}_x\vb*{\mathsf{p}}$ \\
      \stepref{2} & $\mathsf{p}^{(01)} = \bm{\mathsf{F}}_y\phi_y\mathsf{p}^{(0)}$ \\
      \stepref{3} & $\mathsf{p}^{(011)} = \bm{\mathsf{F}}_z\phi_z\mathsf{p}^{(01)}$
    \end{tabular}};
  \node[anchor=north west, font=\footnotesize] at (\Xe-1.6,\by-0.85) {%
    \begin{tabular}{@{}c@{\enspace}l@{}}
      \stepref{4} & $\mathsf{f}^{(011)} = \tilde{\bm{\mathsf{G}}}_\circ^{(011)}\mathsf{p}^{(011)}$ \\
      \stepref{5} & $\mathsf{f}^{(01)} = \tfrac12\ab(\bm{\mathsf{F}}_z^{-1}\mathsf{f}^{(010)} + \phi_z^{-1}\bm{\mathsf{F}}_z^{-1}\mathsf{f}^{(011)})$ \\
      \stepref{6} & $\mathsf{f}^{(0)} = \tfrac12\ab(\bm{\mathsf{F}}_y^{-1}\mathsf{f}^{(00)} + \phi_y^{-1}\bm{\mathsf{F}}_y^{-1}\mathsf{f}^{(01)})$ \\
      & $\vb*{\mathsf{f}} = \tfrac12\ab(\bm{\mathsf{F}}_x^{-1}\mathsf{f}^{(0)} + \phi_x^{-1}\bm{\mathsf{F}}_x^{-1}\mathsf{f}^{(1)})$
    \end{tabular}};
\end{tikzpicture}
